%Version 3.1 December 2024
% See section 11 of the User Manual for version history
%
%%%%%%%%%%%%%%%%%%%%%%%%%%%%%%%%%%%%%%%%%%%%%%%%%%%%%%%%%%%%%%%%%%%%%%
%%                                                                 %%
%% Please do not use \input{...} to include other tex files.       %%
%% Submit your LaTeX manuscript as one .tex document.              %%
%%                                                                 %%
%% All additional figures and files should be attached             %%
%% separately and not embedded in the \TeX\ document itself.       %%
%%                                                                 %%
%%%%%%%%%%%%%%%%%%%%%%%%%%%%%%%%%%%%%%%%%%%%%%%%%%%%%%%%%%%%%%%%%%%%%

%%\documentclass[referee,sn-basic]{sn-jnl}% referee option is meant for double line spacing

%%=======================================================%%
%% to print line numbers in the margin use lineno option %%
%%=======================================================%%

%%\documentclass[lineno,pdflatex,sn-basic]{sn-jnl}% Basic Springer Nature Reference Style/Chemistry Reference Style

%%=========================================================================================%%
%% the documentclass is set to pdflatex as default. You can delete it if not appropriate.  %%
%%=========================================================================================%%

%%\documentclass[sn-basic]{sn-jnl}% Basic Springer Nature Reference Style/Chemistry Reference Style

%%Note: the following reference styles support Namedate and Numbered referencing. By default the style follows the most common style. To switch between the options you can add or remove “Numbered” in the optional parenthesis. 
%%The option is available for: sn-basic.bst, sn-chicago.bst%  
 
%%\documentclass[pdflatex,sn-nature]{sn-jnl}% Style for submissions to Nature Portfolio journals
%%\documentclass[pdflatex,sn-basic]{sn-jnl}% Basic Springer Nature Reference Style/Chemistry Reference Style
\documentclass[pdflatex,iicol,sn-mathphys-num]{sn-jnl}% Math and Physical Sciences Numbered Reference Style
%%\documentclass[pdflatex,sn-mathphys-ay]{sn-jnl}% Math and Physical Sciences Author Year Reference Style
%%\documentclass[pdflatex,sn-aps]{sn-jnl}% American Physical Society (APS) Reference Style
%%\documentclass[pdflatex,sn-vancouver-num]{sn-jnl}% Vancouver Numbered Reference Style
%%\documentclass[pdflatex,sn-vancouver-ay]{sn-jnl}% Vancouver Author Year Reference Style
%%\documentclass[pdflatex,sn-apa]{sn-jnl}% APA Reference Style
%%\documentclass[pdflatex,sn-chicago]{sn-jnl}% Chicago-based Humanities Reference Style

%%%% Standard Packages
%%<additional latex packages if required can be included here>

\usepackage{graphicx}%
\usepackage{multirow}%
\usepackage{amsmath,amssymb,amsfonts}%
\usepackage{amsthm}%
\usepackage{mathrsfs}%
\usepackage[title]{appendix}%
\usepackage{xcolor}%
\usepackage{textcomp}%
\usepackage{manyfoot}%
\usepackage{booktabs}%
\usepackage{algorithm}%
\usepackage{algorithmicx}%
\usepackage{algpseudocode}%
\usepackage{listings}%
\usepackage[all]{nowidow}
\usepackage[none]{hyphenat}

%%%%

%%%%%=============================================================================%%%%
%%%%  Remarks: This template is provided to aid authors with the preparation
%%%%  of original research articles intended for submission to journals published 
%%%%  by Springer Nature. The guidance has been prepared in partnership with 
%%%%  production teams to conform to Springer Nature technical requirements. 
%%%%  Editorial and presentation requirements differ among journal portfolios and 
%%%%  research disciplines. You may find sections in this template are irrelevant 
%%%%  to your work and are empowered to omit any such section if allowed by the 
%%%%  journal you intend to submit to. The submission guidelines and policies 
%%%%  of the journal take precedence. A detailed User Manual is available in the 
%%%%  template package for technical guidance.
%%%%%=============================================================================%%%%

%% as per the requirement new theorem styles can be included as shown below
\theoremstyle{thmstyleone}%
\newtheorem{theorem}{Theorem}%  meant for continuous numbers
%%\newtheorem{theorem}{Theorem}[section]% meant for sectionwise numbers
%% optional argument [theorem] produces theorem numbering sequence instead of independent numbers for Proposition
\newtheorem{proposition}[theorem]{Proposition}% 
%%\newtheorem{proposition}{Proposition}% to get separate numbers for theorem and proposition etc.

\theoremstyle{thmstyletwo}%
\newtheorem{example}{Example}%
\newtheorem{remark}{Remark}%

\theoremstyle{thmstylethree}%
\newtheorem{definition}{Definition}%

\raggedbottom
%%\unnumbered% uncomment this for unnumbered level heads

\begin{document}

\title[Automatic Prosodic Audio Description Modeling]{Automatic Prosodic Audio Description Modeling}

%%=============================================================%%
%% GivenName    -> \fnm{Joergen W.}
%% Particle     -> \spfx{van der} -> surname prefix
%% FamilyName   -> \sur{Ploeg}
%% Suffix       -> \sfx{IV}
%% \author*[1,2]{\fnm{Joergen W.} \spfx{van der} \sur{Ploeg} 
%%  \sfx{IV}}\email{iauthor@gmail.com}
%%=============================================================%%

\author*[1]{\fnm{Christian} \sur{Quintero} %\orcid{0000-0001-5763-2373}%
}\email{christian.quintero@unimilitar.edu.co}
\author[1]{\fnm{Alexander} \sur{Rozo-Torres} %\orcid{0000-0001-9343-2264}%
}\email{est.omar.rozo@unimilitar.edu.co} 
\equalcont{These authors contributed equally to this work.}
\author[1]{\fnm{Cristian} \sur{Plazas} %\orcid{0009-0008-2633-9607}%
}\email{est.cristian.plazas@unimilitar.edu.co} 
\equalcont{These authors contributed equally to this work.}



\affil*[1]{\orgname{Universidad Militar Nueva Granada},
           \orgaddress{\city{Bogotá},
                       \country{Colombia}}}



%%==================================%%
%% Sample for unstructured abstract %%
%%==================================%%

\abstract{
\noindent \textbf{Purpose:} Current audio description systems struggle to achieve semantic accuracy and emotional expressiveness simultaneously. This work proposes a unified framework for prosodic audio description modeling integrated with semantic narrative generation, aimed at enhancing audiovisual accessibility for visually impaired people, leveraging evidence that these users exhibit heightened sensitivity to prosodic nuance.

\noindent \textbf{Methods:} Descriptions of ten representative MSVD video scenes were analyzed, comparing original annotations against narratives generated through semantic and attention-based refinement. Prosodic features—fundamental frequency, intensity, speech rate, and pause duration—were modeled using 16 emotional renditions from three professional narrators. Acoustic features were extracted via OpenSMILE (eGeMAPSv02), and emotional target profiles were computed across narrators. Two synthesis strategies were evaluated: prompt-based emotional TTS and reference-audio voice cloning (XTTS v2, Spanish-F5), conditioned on target prosodic profiles derived from Plutchik's emotion wheel.

\noindent \textbf{Results:} Analysis of 60 synthesized audio descriptions reveals that reference-based synthesis achieved greater congruence with target profiles than prompt-based alternatives, particularly for high-arousal emotions. High-arousal emotions exhibited elevated fundamental frequency, higher intensity, and accelerated speech rate, whereas low-arousal emotions displayed descending frequency, reduced intensity, and longer pauses. Patterns remained consistent across narrators, indicating robust emotion-level prosodic markers.

\noindent \textbf{Conclusion:} The framework appears promising for enhancing descriptive accuracy and emotional expressiveness, potentially transcending the monotony of conventional audio description systems. However, perceptual validation with visually impaired users is essential to confirm actual accessibility improvements and effectiveness.
}


\keywords{audio description, prosodic modeling, semantic narrative, emotional speech synthesis, audiovisual accessibility}

%%\pacs[JEL Classification]{D8, H51}

%%\pacs[MSC Classification]{35A01, 65L10, 65L12, 65L20, 65L70}

\maketitle

\section{Introduction}\label{sec1}

Equitable access to audiovisual content remains a critical challenge for people who are blind or low vision (BLV). Accurate, well-adapted Audio Descriptions (AD) are still scarce, and existing approaches are often insensitive to the narrative complexity and emotional texture of video~\cite{Sharma2023,Islam2021}. In recent years, automatic AD based on generative models has emerged as a promising solution, as joint advances in computer vision and natural language processing (NLP) make it possible to produce textual descriptions that render visual content in increasingly natural and inclusive ways~\cite{Shen2024a,Yuan2025}. Nevertheless, systems originally designed for static images struggle with the temporal and semantic richness of video and, in particular, with the fundamental challenge of conveying emotion through synthesized speech~\cite{Taha2024,Cusworth2024}.

AD is commonly defined as an accessibility technique that translates the dialogue-free visual elements of a scene into verbal language, thereby enabling BLV people to understand and enjoy audiovisual content~\cite{Shen2024a,Yucel2025}. Within this framework, prosody, understood as the suprasegmental features of speech that shape \textit{how} something is said beyond the literal content of the words, plays a fundamental role: it not only contributes to speech intelligibility but also conveys communicative intent, emotion, and narrative meaning~\cite{Iavarone2024,Gu2024}. In particular, intonation can emphasize moments of tension or calm, while rhythm and pauses mark narrative transitions or highlight critical details, thereby complementing linguistic content and enriching the listening experience~\cite{Nussbaum2023,Iavarone2024}.

In automatic AD, \textit{semantic narrative}, referred to in the computer vision literature as \textit{dense video captioning}, denotes the ability to generate textual descriptions that do not merely enumerate objects or actions but build a coherent, contextualized account adapted to the temporal dynamics of the scene~\cite{Zhou2018,Li2018}. This approach goes beyond simple visual identification, incorporating information about spatial relationships, event sequences, and contextual nuances that enrich overall comprehension. By drawing on transformer-based architectures and convolutional neural networks, events can be described without manual segmentation, yielding detailed, temporally coherent sequences~\cite{Zhu2019,Yang2023}.

Building on these capabilities, multimodal fusion, which combines visual, auditory, and textual information, has been shown to further reinforce descriptive coherence and contextual richness~\cite{Zhao2020}. However, even when descriptions are semantically rich, their delivery through conventional text-to-speech (TTS) models tends to exhibit limited prosodic expressiveness, producing monotonous narrations that may undermine listener immersion and weaken emotional nuance~\cite{Taha2024,Qi2024}.

Indeed, current TTS systems show notable limitations in generating natural, expressive prosody, which can negatively affect perceived content quality~\cite{Oh2024,Cornille2022}. In AD specifically, this lack of prosodic variability may lead to a monotonous experience, diminishing the conveyance of emotional nuance and, in turn, potentially hindering comprehension and enjoyment of BLV people~\cite{Liu2023,Ben-David2024}. To address this gap, emotional TTS (ETTS) models have sought to reproduce human prosodic patterns from reference audio, textual prompts, or latent representations, achieving more natural and adaptive results~\cite{Bott2024,Cho2024}.

From a theoretical standpoint, Plutchik's wheel~\cite{Plutchik2001,Plutchik1989} and Russell's circumplex model~\cite{Russell1980} constitute the foundational frameworks most widely used to map emotional expression in speech. These models provide solid conceptual scaffolding for understanding how emotions are structured across the valence-arousal space and, in turn, offer useful tools for designing ETTS systems~\cite{Zheng2025,Zhou2023}.

In parallel, cognitive neuroscience research indicates that individuals with late-onset blindness exhibit heightened sensitivity to the emotional nuances of the voice, with enhanced auditory cortex recruitment for affective speech~\cite{Ben-David2024,Liu2024a}. Findings on affective authenticity further suggest that even subtle prosodic artificiality is readily perceived by this population, which may directly harm immersion~\cite{Sarzedas2024,Nussbaum2023}. Synthesis intended for AD must therefore strive to balance acoustic naturalness with emotional clarity.

Against this backdrop, the present work proposes a unified framework that integrates attention mechanisms, in order to refine the semantic narrative, with emotional prosodic modeling, in order to deliver it expressively. This approach seeks to produce detailed, contextually coherent descriptions while also enabling an emotionally richer listening experience, with the broader goal of advancing digital inclusion for BLV people.

The main motivation stems from a gap that persists despite significant progress in both video captioning and emotional TTS: to the best of our knowledge, no existing automatic AD system jointly addresses semantic richness and prosodic expressiveness within a unified, emotionally grounded pipeline. Current approaches either focus on generating accurate descriptions without considering how they are delivered, or apply generic TTS synthesis that strips the narration of the affective cues that are particularly salient to BLV listeners. Bridging this gap is not merely a technical concern; it has direct implications for the lived experience of a population for whom speech is a primary channel of audiovisual access. By grounding prosodic modulation in a well-established psychological model of emotion and validating it through large-scale acoustic analysis, this work contributes both a practical system and a reproducible evaluation methodology on which future research can build.

This paper is organized as follows. Section~\ref{sec:related} reviews related work on dense video captioning, attention mechanisms, emotional prosodic modeling, and auditory perception in BLV users. Section~\ref{sec:framework} describes the proposed framework and its four stages. Section~\ref{sec:setup} details the experimental setup, including the corpus, synthesis strategies, and acoustic evaluation protocol. Section~\ref{sec:results} presents the results across the three evaluation dimensions. Section~\ref{sec:discussion} discusses the findings and their implications. Section~\ref{sec:conclusions} draws the conclusions, and Section~\ref{sec:future} outlines directions for future work.

\section{Related Work}
\label{sec:related}
% ----

Automatic AD sits at the intersection of computer vision, NLP, and speech synthesis, three fields that have each advanced rapidly yet whose integration remains an open research challenge. This section reviews the conceptual axes most relevant to the present work: dense video captioning and semantic narrative generation, attention mechanisms and multimodal fusion, and emotional prosodic modeling for TTS.

The first axis, semantic narrative, referred to in the literature as \textit{dense video captioning}, targets the generation of detailed, coherent descriptions of multiple events that coexist within a single video. Foundational architectures for jointly localizing and describing events without manual segmentation were established using transformer-based and proposal-driven approaches~\cite{Zhou2018,Li2018}. Densely connected recurrent networks that retain long-range context have further improved action-sequence accuracy~\cite{Zhu2019}, while large-scale visual-language models pretrained on narrated videos excel at temporally grounded captioning~\cite{Yang2023}. More recent work demonstrates that fine-grained temporal reasoning substantially improves long-video understanding~\cite{Yuan2025}, and that controllable caption length can be addressed via ordinal embeddings~\cite{Nitta2024}. Zero-shot approaches that generate detailed descriptions without task-specific training ease scalability to low-resource domains~\cite{Xie2024}, and multimodal memory networks that integrate visual, textual, and auditory cues produce stylized, coherent narratives~\cite{Zhao2020}. Collectively, these contributions consolidate a robust framework for descriptive generation; none of them, however, explicitly addresses the prosodic dimension.

The second axis, attention and multimodal fusion, concerns how attention mechanisms have become a cornerstone for dynamically filtering the most relevant information within an audiovisual stream. CNN-LSTM architectures with hierarchical attention notably reduce descriptive noise and improve temporal coherence~\cite{Huang2022}, whereas cross-modal attention that aligns visual and textual representations in a shared semantic space reinforces informational density and enables more natural narrative transitions~\cite{Li2021}. In addition, combining visual and auditory features increases contextual richness in automatic AD tools~\cite{Shen2024a}, and attention-driven pipelines have been shown to generalize effectively to specialized visualization domains~\cite{Smits2024}. The resulting textual output is therefore not only accurate in content but also pragmatically adapted to the communicative situation.

The third axis concerns emotional prosodic modeling. Although semantic accuracy has improved substantially, prosodic control remains comparatively underexplored in TTS as applied to AD. Early work in this direction has shown that semi-supervised approaches can generate emotionally expressive speech without explicit emotion annotations~\cite{Cai2020}, while multilevel prosodic control within neural TTS architectures enables fine-grained adjustments down to the phoneme level~\cite{Cornille2022}. In parallel, sequence-to-sequence attention has been used to modulate emotion duration~\cite{Choi2021}, and non-parallel voice conversion via cycle-consistent adversarial networks avoids the over-smoothing of traditional statistical models~\cite{Kaneko2018}.

Latent-space methods constitute a complementary line of research within this axis. In particular, Global Style Tokens (GST)~\cite{Wu2019,Wang2018} and Variational Autoencoders (VAE)~\cite{Zhang2018} enable emotion transfer with limited labeled data. More recently, diffusion-based latent prosody generation with adversarial training has achieved highly expressive synthesis~\cite{Oh2024}, multilingual emotion transfer with minimal supervision has extended these results across languages and speakers~\cite{Zhu2023}, and fine-grained prosodic control through natural-language prompts offers a flexible alternative to reference audio~\cite{Bott2024}. Emotional voice conversion using deep bidirectional architectures with Mel-frequency cepstral coefficients (MFCC) features~\cite{Cao2024} and contextual generative models~\cite{Porco2023} further illustrate the breadth of available strategies. Despite this progress, the application of ETTS to automatic AD remains incipient.

Finally, BLV people, particularly those with late-onset blindness, show heightened sensitivity to the emotional nuances of the voice, with enhanced auditory cortex specialization for affective speech processing~\cite{Ben-David2024,Liu2024a}. Poorly calibrated prosodic artificiality is readily perceived by this population, which can directly undermine immersion~\cite{Sarzedas2024,Nussbaum2023}. Perceptual metrics for quantifying emotional similarity in speech~\cite{Harvill2023} and documented temporal and spectral markers of emotional speech~\cite{Kurazhova2024} are directly relevant to designing evaluation criteria for AD systems. Moreover, research on voice-based interfaces for BLV people in navigation contexts confirms that prosodic quality is a key driver of user experience~\cite{Yucel2025}. Together, these findings motivate the design of AD synthesis that balances acoustic naturalness with emotional clarity, adapted to the perceptual expectations of BLV people.

\section{Proposed Framework}
\label{sec:framework}
% ----

This paper proposes a modular framework that integrates image analysis, contextual semantic extraction, and automated prosodic characterization in order to optimize the generation of emotionally expressive AD. The overall goal is twofold: to improve accessibility for BLV users and to enrich the narrative experience by aligning the synthesized voice with the affective content of the scene. Fig.~\ref{fig:esquema} illustrates the architecture, which is organized into four interrelated stages that progress from low-level audiovisual segmentation to high-level prosodic rendering.

\begin{figure*}[htbp]
    \centering
    \includegraphics[width=.95\textwidth]{Images/AD_prosodico.pdf}
    \caption{Overview of the proposed Automatic Prosodic Audio Description Modeling. The pipeline applies attention mechanisms over key scene elements, generates a semantic description, and adjusts the prosodic features of speech synthesis according to an emotional scheme grounded in Plutchik's model~\cite{Plutchik2001}.}
    \label{fig:esquema}
\end{figure*}

\subsection{Attention Mechanisms}
\label{subsec:attention}

In the first stage, \textbf{attention mechanisms} are applied to the audiovisual scene to prioritize the information most relevant to scene comprehension, thereby avoiding descriptive overload. Seven complementary mechanisms were defined, each targeting a distinct visual-narrative dimension:

\begin{itemize}
    \item \textbf{Actions:} movements, gestures, or behaviors executed by characters that drive the narrative or communicate intentions and emotions. From a semantic perspective, actions are minimal units of meaning that articulate the logic of events; their mention in AD must emphasize what is relevant for narrative and emotional comprehension.

    \item \textbf{Special effects:} visual, auditory, or digital elements that introduce a spectacular, symbolic, or fantastical dimension into the scene. Their semantic function may reinforce emotions, emphasize actions, or create transitions; in AD, they should be interpreted in terms of emotional and narrative impact rather than technical description.

    \item \textbf{Environment and atmosphere:} the physical setting combined with its sensory and affective charge, which together configure the emotional and symbolic backdrop of the narrative action. This category encompasses both material elements of the environment and its expressive qualities, and its translation in AD should prioritize sensory aspects that shape the listener's emotional experience.

    \item \textbf{Facial expressions:} muscular configurations of the face that communicate emotions, intentions, or internal reactions. They function as iconic signs connecting the character's interiority with the viewer, and AD must translate them into terms that evoke emotional intensity and quality with precision.

    \item \textbf{Interactions:} dynamic relationships between characters or between characters and objects or environments, revealing affective, social, or symbolic bonds. These are key to understanding roles, hierarchies, conflicts, or alliances within the narrative universe.

    \item \textbf{Characters:} narrative entities with physical, psychological, social, and emotional traits whose presence carries semantic weight in the construction of the story. Their description should prioritize attributes relevant to the plot and aspects that mark their emotional evolution.

    \item \textbf{Body language:} physical configurations and non-verbal movements that express emotional states, attitudinal dispositions, or interpersonal relations. In AD, they must be translated into verbal images that allow the listener to infer emotional state, intention, or narrative tension.
\end{itemize}

Taken together, these seven dimensions define the descriptive scope that subsequent stages aim to verbalize, and they constitute the criteria according to which keyframes and visual cues are subsequently selected.

\subsection{Semantic Description Generation}
\label{subsec:semantic}

Building on the prioritized cues produced by the attention layer, the second stage generates a detailed \textbf{semantic description} of the scene. This description is intended as a coherent, contextualized account that is adapted to the temporal dynamics of the scene and supported by natural language processing (NLP) and computer vision models, in order to ensure alignment between visual content and verbal description. To this end, the semantic description stage is decomposed into three tightly coupled sub-processes, summarized below and depicted in Fig.~\ref{fig:esquema}.

The pipeline begins with a \textbf{reverse diarization} step, in which the audio track is analyzed to detect intervals containing speech from on-screen or off-screen speakers. By inverting the resulting diarization mask, the system isolates the \emph{speaker-free} segments of the video, i.e., the intervals in which an AD track can be inserted without overlapping with original dialogue. This speech-aware partitioning enables a scene-independent yet temporally structured segmentation of the content and provides the temporal anchors on which the remaining stages of the framework operate.

Within each speaker-free segment, semantic content is extracted from the visual stream by means of CLIP (Contrastive Language–Image Pre-training)~\cite{Radford2021}. CLIP jointly trains an image encoder and a text encoder under natural-language supervision so as to predict the correct pairings between images and textual captions, which allows it to acquire transferable visual concepts without task-specific labels. In our framework, CLIP is used as a multimodal embedding space in which candidate frames are ranked against textual prompts derived from the seven attention dimensions described in Section~\ref{subsec:attention}. Frames that maximize the image–text similarity for the most salient prompts are retained as \textbf{keyframes}, yielding a compact set of representative views per segment. This selection step concentrates the descriptive effort on the visual evidence that is most likely to carry narrative and emotional weight.

From the selected keyframes, detailed descriptions are then produced using generative language models equipped with attention-based architectures, which weight input elements according to their relevance and thus focus the description on the most informative visual cues. We refer to the output of this process as the \textbf{semantic narrative}, since it is centered on identifying and verbalizing the objects, actions, and contexts that are most pertinent to each segment, rather than on exhaustively enumerating visual details. The resulting descriptions aim to be precise, temporally coherent, and faithful to the underlying scene, while remaining concise enough to fit within the available speaker-free intervals.

In parallel with the textual narrative, a set of \textbf{prosodic features}—namely intensity, fundamental frequency, rhythm, and pause duration—is associated with each segment as a function of its inferred emotional intention. This prosodic component is intended to enrich the auditory experience by ensuring that the narration is not only semantically accurate but also emotionally aligned with the depicted scene. The mapping from emotional intention to prosodic parameters is formalized in the following stage.

\subsection{Emotional Characterization}
\label{subsec:emotional}

In the third stage, the system incorporates the \textbf{emotional dimension} by identifying the predominant emotion in the scene according to Plutchik's model~\cite{Plutchik2001,Plutchik1989}. Sixteen emotional interpretations are considered, each linked to a primary emotion of the wheel and associated with a specific prosodic profile defined by fundamental frequency, intensity, rhythm, and pause duration. This mapping provides the bridge between the semantic narrative produced in the previous stage and the prosodic rendering performed in the next one.

Each generated description was associated with a predominant emotion (Joy, Trust, Fear, Surprise, Sadness, Disgust, Anger, and Anticipation) according to Plutchik's model~\cite{Plutchik2001}. A reference corpus of 16 emotional renditions was recorded by three narrators (Woman01, Man01, and Man02), covering 16 emotion categories across high- and low-intensity variants, which produced 48 audio references. Four prosodic parameters were extracted from these recordings using OpenSMILE: fundamental frequency (F0), intensity (energy-based RMS), speech rate (syllables per second), and pause duration (inter-phrase silence). Per-emotion average values were computed and used as target profiles to parameterize synthesis. Fig.~\ref{fig:radar_locutores} illustrates these prosodic profiles for each of the 16 emotional renditions, grouped by primary emotion and intensity level, with the individual narrator profiles and their mean overlay. By processing keyframes and generated descriptions together, secondary emotions were identified and weighted by relevance, ensuring that the AD is emotionally congruent with the visual content.

\begin{figure}[htbp]
    \centering
    \includegraphics[width=\linewidth]{Images/EmocionesLocutores.pdf}
    \caption{Prosodic profiles of the 16 emotional interpretations recorded by the three narrators, with the group mean represented by a dashed gray line. Each radar displays the four prosodic dimensions normalized on a scale of [1,5]. The top row corresponds to high-intensity variants and the bottom row to low-intensity variants of each primary emotion on Plutchik's wheel.}
    \label{fig:radar_locutores}
\end{figure}

\subsection{Prosodic Audio Description Synthesis}
\label{subsec:prosodic}

In the final stage, the \textbf{prosodic AD} adjusts its prosodic parameters to reflect the detected emotion, with the aim of producing AD that enriches the listening experience and facilitates user immersion. The system is grounded in the valence–arousal space shared by Plutchik's and Russell's frameworks~\cite{Plutchik2001,Russell1980}: high-arousal emotions such as \textit{joy} or \textit{anger} tend to exhibit elevated F0, higher intensity, and faster speech rate, whereas low-arousal emotions such as \textit{sadness} or \textit{fear} are typically characterized by descending F0, lower intensity, and longer pauses. This correspondence is visualized in Fig.~\ref{fig:emociones}.

\begin{figure}[htbp]
    \centering
    \includegraphics[width=0.95\linewidth]{Images/EmocionesBase.pdf}
    \caption{Correspondence between Plutchik's emotion wheel and prosodic parameters of synthesized speech. Each emotion is associated with specific variations in intensity (I), fundamental frequency/pitch (T), speech rate (R), and pause duration (D).}
    \label{fig:emociones}
\end{figure}





Acoustic evaluation relies on OpenSMILE\footnote{OpenSMILE (Open-Source Media Interpretation by Large feature-space Extraction), developed by audEERING GmbH and the Technical University of Munich. \url{https://audeering.github.io/opensmile/}}, an open-source toolkit widely adopted for large-scale, low-level acoustic feature extraction from audio streams. OpenSMILE enables automated computation of F0 statistics, energy-based intensity measures, rhythm descriptors, and pause segmentation from arbitrary audio batches via Python, which makes it well suited for the volume of synthesized AD produced in this study.

Because the four descriptors are expressed in heterogeneous units, each feature was standardized via a per-feature z-score and then linearly rescaled to a common $[1,5]$ interval. This normalization yields dimensionless, directly comparable prosodic profiles across emotions, narrators, and synthesis models, and underlies the radar representations.

For each synthesized AD, the normalized four-dimensional prosodic vector was compared against the per-emotion target profile derived from the human reference corpus. Alignment was quantified as the mean absolute deviation between the generated and target profiles across the four dimensions, and profiles were aggregated by emotion, by video, and by narrator for comparative analysis.

% ----


% -------------------------------------------------------
\section{Methodology and Experimental Setup}
\label{sec:setup}
% -------------------------------------------------------

To develop an emotionally expressive AD pipeline, we designed a systematic workflow integrating computer vision, NLP, and TTS synthesis. Fig.~\ref{fig:metodología} presents the six sequential phases: scene selection and segmentation, semantic description generation, emotional prosodic analysis and modeling, model comparison, speech synthesis, and acoustic evaluation.

\begin{figure*}[ht]
    \centering
    \includegraphics[width=1\linewidth]{Images/Metodologia.pdf}
    \caption{Flow diagram of the proposed framework, detailing the six sequential phases: scene selection, description generation, prosodic modeling, model comparison, speech synthesis, and acoustic evaluation.}
    \label{fig:metodología}
\end{figure*}

\textbf{Scene Selection and Segmentation:} We used the Microsoft Research Video Description Corpus (MSVD)\footnote{MSVD, Microsoft Research. Available at: \url{https://github.com/microsoft/Microsoft-Research-Video-Description-Corpus}}~\cite{Yang2023}, comprising approximately 1,800 short video clips each accompanied by multiple manually produced descriptions in several languages~\cite{Islam2021}. The corpus is partitioned into training (1,200 videos), validation (100), and test (670) sets. All videos have their audio muted; clips with subtitles or overlaid text are excluded. Video duration typically ranges from 10 to 25 seconds, with most depicting a single activity. MSVD is widely used in AD, human-robot interaction, video surveillance, and automatic scene understanding~\cite{Sharma2023}.

Ten representative scenes were selected (Fig.~\ref{fig:Contexto_videos}), prioritizing high emotional load and visual diversity. Keyframes were extracted using Contrastive Language-Image Pretraining (CLIP)~\cite{Radford2021} with adaptive segmentation dynamically adjusted to each video's duration and frame rate; no original audio was present in the clips.

\begin{figure}[htbp]
    \centering
    \includegraphics[width=0.4\textwidth]{Images/Contexto videos.pdf}
    \caption{Videos used in the study with their annotations from the MSVD corpus.}
    \label{fig:Contexto_videos}
\end{figure}

\textbf{Semantic Description Generation:} From the extracted keyframes, automatic descriptions were generated using GPT-4o\footnote{GPT-4o, OpenAI. \url{https://openai.com/index/hello-gpt-4o/}}, with a character limit proportional to video duration (30 characters per second). Advanced attention mechanisms were incorporated to improve coherence and detail, prioritizing the most relevant scene elements~\cite{Huang2022,Li2021}.

\textbf{Speech Synthesis Strategies:} To evaluate prosodic expressiveness under different conditioning paradigms, two complementary synthesis strategies were applied. The first relies solely on textual prompts, while the second leverages human-recorded reference audio to achieve finer prosodic control.

\begin{enumerate}
    \item \textbf{Prompt-based generation:} Explicit textual instructions were used to induce target emotions in TTS models supporting prompt control.
    \item \textbf{Reference-audio generation (primary approach):} Recordings from the emotional reference corpus were used as acoustic conditioning inputs for voice-cloning models to replicate the target prosodic patterns.
\end{enumerate}

For the reference-audio strategy, two open-source voice-cloning models were selected based on their capacity to reproduce prosodic characteristics from a conditioning recording:

\begin{itemize}
    \item \textbf{XTTS v2}\footnote{XTTS v2, Coqui TTS (open source). \url{https://github.com/coqui-ai/TTS}}: a zero-shot multilingual voice-cloning model recognized for faithfully reproducing the pitch, rhythm, and intensity of reference recordings.
    \item \textbf{Spanish-F5}\footnote{Spanish-F5, open-source implementation of F5-TTS for Spanish. \url{https://huggingface.co/jpgallegoar/Spanish-F5}}: a Spanish-language implementation of the F5-TTS architecture trained on 218 hours of audio spanning Latin American and Peninsular Spanish accents, used for its capacity to replicate reference-voice prosodic characteristics.
\end{itemize}

Semantic content was held constant across all outputs; only prosodic parameters varied according to each tool's configuration. Each model produced one AD per video, except the reference-audio models (XTTS v2 and Spanish-F5), each of which used three reference clips (one per narrator), yielding three AD versions per video. The synthesis process produced 60 audio files: 2~models~$\times$~10~videos~$\times$~3~narrators.

\textbf{Acoustic Evaluation:} Acoustic evaluation focused on the expressiveness, naturalness, and emotional appropriateness of the synthesized AD. All prosodic features were extracted automatically using OpenSMILE, configured with the \textit{eGeMAPSv02} feature set at the Functionals level, from which the four target parameters were derived. Prosodic parameters were extracted and normalized for analysis:

\textbf{Intensity (I):} mean perceived loudness, conveying emphasis, force, or emotional softness (\texttt{loudness\_sma3\_amean}).

\textbf{Pitch / F0 (T):} mean fundamental frequency on a semitone scale (ref.\ 27.5~Hz), associated with intonation and emotional perception (\texttt{F0semitoneFrom27.5Hz\_sma3nz\_amean}).

\textbf{Rhythm (R):} voiced segments per second, a tempo proxy indicating dynamism or deliberateness (\texttt{VoicedSegmentsPerSec}).

\textbf{Pause duration (D):} mean unvoiced segment length, a pause proxy affecting rhythm and speech segmentation (\texttt{MeanUnvoicedSegmentLength}).

The extended Geneva Minimalistic Acoustic Parameter Set (eGeMAPS) was selected because it is a compact, expert-curated standard explicitly designed for affective and clinical voice research, which improves comparability and reproducibility relative to brute-force feature sets. OpenSMILE was chosen for its robustness in batch processing, its native Python interface, and its established use in speech emotion recognition research, making it well suited to the corpus of synthesized AD produced in this study.

% -------------------------------------------------------
\section{Results}
\label{sec:results}
% -------------------------------------------------------

Building on the experimental protocol described in Section~\ref{sec:setup}, the evaluation was conducted on the 10 MSVD videos, each associated with one of eight primary emotions from Plutchik's wheel. The 60 audio files ($10~\text{videos} \times 3~\text{narrators} \times 2~\text{models}$) were analyzed with OpenSMILE under the eGeMAPSv02 Functionals configuration. Analysis was organized around three dimensions: semantic coherence, prosodic expressiveness, and emotional alignment.

\textbf{Model Comparison for Description Generation:} Eight generative AI models for video description were comparatively evaluated using the official MSVD reference annotations for 10 videos in Spanish, chosen to ensure diversity of scenes and contexts. Semantic similarity was assessed using the CIDEr, BLEURT, and ROUGE metrics. Based on these results, GPT-4o (used in our model) was selected as the basis for the description, and a systematic AD generation protocol was established that encompasses narrative reuse, emotional TTS synthesis, and comparative prosodic analysis, as shown in Fig.~\ref{fig:model_descriptions_metrics}.

\begin{figure*}[htbp]
    \centering
    \includegraphics[width=0.95\linewidth]{Images/benchmark_metrics.png}
    \caption{Comparison of average semantic similarity scores (CIDEr, BLEURT, and ROUGE) across eight generative AI models evaluated against the Spanish reference annotations. GPT-4o (Our model) achieved the highest overall performance.}
    \label{fig:model_descriptions_metrics}
\end{figure*}

\textbf{Semantic coherence:} BLEU, CIDEr, and ROUGE metrics indicated that GPT-4o obtained the highest scores against the Spanish reference annotations. Its descriptions were consistently rated as more accurate, detailed, and contextually appropriate than those of the seven alternative models.

\textbf{Prosodic expressiveness:} The reference-audio-based models (XTTS v2 and Spanish-F5) achieved greater average variability in pitch, intensity, and speech rate, closely approximating the profiles defined in the emotional reference corpus. By contrast, the prompt-controlled models, although enabling explicit emotion specification, exhibited less pronounced prosodic modulation.

\textbf{Emotional alignment:} The correspondence between the target emotion and the generated prosodic profiles was confirmed by inter-narrator analysis. Differences between narrators emerged primarily on intensity and pitch, reflecting individual vocal styles. By contrast, rhythm remained largely consistent across narrators, confirming that tempo cues are robustly shared as markers of emotional category regardless of speaker.

\subsection{Prosody Assignment}

\noindent For the video descriptions, a generative artificial intelligence approach was adopted in which the keyframes of each video were analyzed alongside the description generated by our model. Table~\ref{tab:evaluacionemoc} presents this analysis, where the most suitable prosodic characteristics for narration were determined, considering primary emotions, intensity, pitch, pause duration, and rhythm.

\begin{table*}[htbp]
    \centering
    \caption{Prosodic profile assignment per video based on JSON analysis (10 MSVD clips).}
    \label{tab:evaluacionemoc}
    \renewcommand{\arraystretch}{1}
    \setlength{\tabcolsep}{4pt}
    \small
    \begin{tabular}{p{2.5cm}p{2cm}p{1.3cm}p{2.6cm}p{2.4cm}p{3.4cm}}
        \toprule
        \textbf{VideoID} & \textbf{Primary Emotion / Comp.} & \textbf{Intensity} & \textbf{Pitch} & \textbf{Duration} & \textbf{Rhythm} \\
        \midrule
        \textbf{XZ0Ch\_ClBfo} & Joy/ Ecstasy & High & High, wide variation & Short, with exclamations & Fast, energetic, dynamic \\
        \textbf{YEmS8ZQIhTA} & Anticipation/ Vigilance & High & High, at times tense & Moderate, with pauses & Agile, fast, alert \\
        \textbf{YersIyzsOpc} & Trust/ Acceptance & Low & Medium, warm, low variability & Moderate & Steady, fluid, no abrupt changes \\
        \textbf{Z3XarHZ3HpY} & Trust/ Admiration & High & High, slightly rising & Moderate & Emphatic, heightened, not abrupt \\
        \textbf{ZdlG8fjGJlo} & Trust/ Acceptance & Low & Medium, warm, low variability & Moderate & Steady, fluid, no abrupt changes \\
        \textbf{yAD\_TS5L2d4} & Anticipation/ Interest & Low & High, slightly rising & Variable & Fluid, agile, sustained attention \\
        \textbf{z0Si1XxMibg} & Trust/ Acceptance & Low & Medium, warm, low variability & Moderate & Steady, fluid, no abrupt changes \\
        \textbf{zTn-nNj5Bng} & Trust/ Acceptance & Low & Medium, warm, low variability & Moderate & Steady, fluid, no abrupt changes \\
        \textbf{za-9mBZyNfQ} & Joy/ Ecstasy & High & High, wide variation & Short, with exclamations & Fast, energetic, dynamic \\
        \textbf{zpgW7m7\_LZw} & Joy/ Serenity & Low & Low, stable, no large variations & Long, relaxed pauses & Slow, regular, measured \\
        \bottomrule
    \end{tabular}
\end{table*}

Beyond the emotion-level analysis, inter-narrator variability reveals consistent individual vocal strategies: Woman01 shows the highest mean pitch (4.06) and the lowest intensity (1.87), reflecting a lighter, more melodic narration style; Man01 produces the highest mean intensity (3.68) with a lower pitch (2.37), characteristic of a more emphatic delivery; and Man02 exhibits the fastest mean rhythm (3.52), suggesting a more dynamic pacing. Despite these individual differences, emotion-level patterns remain largely congruent across narrators.

\begin{figure*}[htbp]
   \centering
    \includegraphics[width=\linewidth]{Images/Modelos_Evaluacion.pdf}
    \caption{Prosodic profiles across videos and narrators extracted with OpenSMILE's normalized scale. Each radar corresponds to a video with its assigned emotion across the four dimensions: intensity (I), pitch (T), speech rate (R), and pause duration (D).}
    \label{fig:evaluacion}
\end{figure*}




% -------------------------------------------------------
\section{Discussion}
\label{sec:discussion}
% -------------------------------------------------------

The results support the central hypothesis that integrating semantic narrative with emotional prosodic modeling substantially improves the quality and expressiveness of automatic AD. Dynamically adjusting prosody to the predominant scene emotion is a crucial step toward a more immersive, accessible listening experience for BLV people.

Evaluation was based on three criteria: accuracy (correct identification of characters, objects, events, and contextual setting), level of detail (relevant information without redundancy), and narrative clarity (coherence and comprehensibility of the generated descriptions).

A key finding is the differential performance of XTTS~v2 and Spanish-F5 in emotional prosodic alignment. Both models condition synthesis on human-recorded reference audio, but they differ in architecture, training data, and voice-cloning strategy. The analysis reveals that their prosodic profiles diverge most on high-arousal emotions (joy, anger, surprise), where fine-grained F0 variability is the primary discriminating marker, and converge more closely on low-arousal emotions (sadness, serenity), where the dominant cues are longer pauses and descending F0 contours. These differences are consistent with findings from the ETTS literature showing that reference-audio conditioning is sensitive to the specific acoustic characteristics of the training protocol.

These results reinforce the importance of an emotionally annotated reference corpus and of synthesis strategies that enable precise prosodic profile transfer. The proposed methodology also demonstrates that it is possible to maintain semantic coherence and contextual appropriateness without sacrificing emotional expressiveness, an outcome directly relevant to the cognitive and perceptual expectations of BLV listeners.

Challenges remain, however, in reproducing subtle, low-arousal emotions and in maintaining synthesis stability across scenes with rapid emotional transitions. Inter-speaker variability and dependence on reference-audio quality also pose generalization challenges. Moreover, the acoustic evaluation conducted here, while robust and reproducible, captures physical signal properties rather than perceived emotional quality. Complementing it with perceptual studies involving BLV end users would provide the missing subjective validation layer.

On the other side, despite promising results, several limitations must be acknowledged. First, the MSVD corpus consists of short clips (10-25 seconds) depicting single, bounded activities. This controlled scope, while useful for analysis, does not reflect the narrative complexity of longer audiovisual content such as feature films or documentaries, which involve extended emotional arcs and complex scene transitions. The brevity of the clips also produced short descriptions, limiting the analysis of multi-clause narrative nuances.

Second, the study was conducted exclusively in Spanish. This restricts the direct applicability of the findings to other languages with different prosodic characteristics and grammatical structures. Regional and dialectal variation within Spanish was also not systematically controlled.

Third, the prosodic evaluation relied solely on acoustic signal features extracted with OpenSMILE. No perceptual validation with BLV end users was conducted in this study; accordingly, the degree to which the measured prosodic alignment translates into subjectively perceived emotional congruence and accessibility remains an open question.

These limitations point to clear directions for future work: evaluation on longer, narratively complex corpora; cross-lingual and cross-dialectal adaptation; and perceptual validation with BLV user populations.

% -------------------------------------------------------
\section{Conclusions}
\label{sec:conclusions}
% -------------------------------------------------------

\noindent This work proposed a unified framework integrating attention-based semantic narrative generation with emotional prosodic modeling for automatic AD. The results indicate that the proposed approach compares favorably with alternative models in semantic accuracy, descriptive richness, and narrative clarity, and that it advances the state of the art in prosodic expressiveness for BLV accessibility.

Attention mechanisms proved instrumental in generating context-enriched descriptions that accurately identify and prioritize the relevant elements of each scene, going beyond the generic or incomplete descriptions typical of conventional approaches. Building on these descriptions, the emotional prosodic modeling implemented here addresses the persistent limitation of expressive monotony in automatic AD: by dynamically adapting fundamental frequency, intensity, speech rate, and pause duration to the emotional load of each scene, the system produces narrations that are more natural and emotionally aligned with the visual content.

Taken together, these contributions support a broader paradigm shift in automatic AD, moving beyond a focus on semantic accuracy alone toward a dual consideration of both the \textit{what} and the \textit{how} of narration. By grounding prosodic modulation in Plutchik's wheel and validating it through large-scale acoustic analysis, the framework offers both a practical pipeline and a reproducible evaluation methodology. This dual contribution enriches the narrative experience and enables the conveyance of not only visual content but also its affective dimension, contributing to a more immersive and inclusive audiovisual experience for BLV users.

% -------------------------------------------------------
\section{Future Work}
\label{sec:future}
% -------------------------------------------------------

\noindent Despite these encouraging results, the present study is bounded by the constraints discussed previously, which in turn open several avenues for future research. A first line of work concerns the corpus itself: extending the evaluation beyond the short MSVD clips to longer, narratively complex content—such as feature films, documentaries, and serialized productions—would allow the framework to be tested under extended emotional arcs and multi-event scene transitions. Within this setting, the dynamic adjustment of characters-per-second, here fixed at 60, also warrants systematic evaluation for its impact on comprehension and audio-visual synchronization.

A second line concerns linguistic and cultural coverage. Because the present analysis was conducted exclusively in Spanish, cross-lingual and cross-dialectal adaptation is needed to assess how the proposed prosodic profiles transfer to languages with different intonational systems and to regional varieties of Spanish. Related to this, prosodic adaptability in short-scene, rapid-emotion contexts requires further investigation; combining the current framework with advanced multimodal fusion strategies could improve descriptive accuracy and emotional tracking for complex or multi-event scenarios.

A third and most critical line is perceptual validation. Future work should incorporate listening studies with BLV end users to validate the emotional quality, naturalness, and accessibility of the generated AD. Only through empirical user feedback will it be possible to calibrate synthesis parameters to real-world preferences and needs, and to determine the extent to which the measured prosodic alignment translates into subjectively perceived emotional congruence.

Finally, this work lays the groundwork for adaptive models that integrate emotional prosody, contextual narrative, and continual learning, leading to TTS systems capable of automatically adjusting narrative and emotional style according to content type, audience profile, or deployment context. Such systems would advance not only the state of the art in automatic AD, but also the broader goal of democratizing access to audiovisual content through emotionally aware technologies.

% -------------------------------------------------------
\section*{Declarations}
% -------------------------------------------------------

\subsection*{Funding}
This work was supported by Universidad Militar Nueva Granada under project INV-ING-4048.

\subsection*{Competing interests}
The authors declare no competing interests.

\subsection*{Ethics approval and consent to participate}
Not applicable. This study did not involve human or animal subjects requiring ethical approval.

\subsection*{Author contributions}
C.~Quintero conceived and supervised the study, designed the framework, and led the writing of the manuscript. C.~Plazas and A.~Rozo-Torres contributed equally to the implementation of the semantic description and prosodic modeling pipeline, the acoustic feature extraction with OpenSMILE, the synthesis experiments with XTTS~v2 and Spanish-F5, and the analysis of the results. All authors reviewed, discussed, and approved the final version of the manuscript.

\subsection*{Use of generative AI}
The authors used GPT-4o for caption generation as part of the experimental pipeline, and general-purpose GPT models for copy-editing and language polishing of the manuscript. No generative AI tool was listed as an author or used to generate scientific content, results, or conclusions.

\bibliography{sn-bibliography}% common bib file
%% if required, the content of .bbl file can be included here once bbl is generated
%%\input sn-article.bbl

\end{document}